% ============================================================
% Section 1.4: Report Organization
% ============================================================
\section{Report Organization}
\label{sec:report_organization}

This thesis is organized into six main chapters structured logically to document the complete engineering lifecycle of the EOG-based assistive system:

\begin{itemize}
    \item \textbf{Chapter 1 (Introduction):} Outlines the clinical background, problem statement, engineering objectives, thesis structure, and execution timeline.
    \item \textbf{Chapter 2 (Literature Review and Theoretical Background):} Reviews existing literature on biopotential acquisition, HMIs, virtual spellers, and eye-tracking systems while identifying current technical gaps.
    \item \textbf{Chapter 3 (System Architecture and Block Diagram):} Details the high-level architecture, functional block diagrams, operational mode-switching logic, and overall system flowcharts.
    \item \textbf{Chapter 4 (Circuit Simulation and Hardware Implementation):} Presents circuit simulations, mathematical design calculations, component selection (AD620, L298N), and software architecture.
    \item \textbf{Chapter 5 (Results and Discussion):} Evaluates experimental biopotential waveforms, obstacle avoidance response times, virtual speller accuracy, and statistical performance metrics.
    \item \textbf{Chapter 6 (Conclusion and Recommendations):} Synthesizes the core engineering findings, clinical applications, and provides strategic recommendations for future technical developments.
\end{itemize}
